\chapter{Introducere}
\label{ch:introducere}

\section{Contextul și motivația lucrării}
În procesul dezvoltării jocurilor de strategie multiplayer în timp real se întâmpină provocări critice: menținerea unei stări de joc sincronizate și minimizarea latenței. Motivația proiectului \textit{Farm Defense: Harvest Wars} pornește de la absența pe piață a unui titlu competitiv de tipul \textit{Tower Defense} cu gameplay de asimetrie totală (unde un jucător este exclusiv atacator, iar celălalt exclusiv apărător). Soluția propusă combină acest tip de \textit{gameplay} cu o arhitectură tehnică modernă: izolarea proceselor, scalabilitate concurentă și decuplarea UI-ului de un server autoritar, prevenindu-se astfel desincronizările și manipularea datelor în scop malițios, lipsit de fair-play (\textit{cheating}).

\section{Stadiul actual al domeniului (Related Work)}
Analiza industriei evidențiază un compromis constant între mecanicile asimetrice și sistemele de progresie. De exemplu, \textit{Plants vs. Zombies} a standardizat deplasarea liniară pe benzi și gestionarea eficientă a economiei interne (\textit{in-game resource management} pentru plasarea unităților), dar rămâne o experiență strict \textit{single-player} (PvE). Pe de altă parte, \textit{Clash Royale} excelează printr-un \textit{metagame} competitiv atractiv, introducând o progresie persistentă a pachetelor de unități în afara meciurilor, spre deosebire de resetarea specifică jocurilor de tip MOBA (ex. \textit{League of Legends}). Totuși, titlurile PvP populare (inclusiv \textit{Bloons TD Battles 2}) forțează o mecanică \textit{simetrică}, în care ambii jucători atacă și se apără simultan pe aceeași arenă.

Lucrarea de față integrează economia dinamică pe benzi și progresia persistentă a unităților într-un mediu PvP competitiv cu asimetrie totală. Pentru a susține tehnic acest sistem de \textit{metagame} și a preveni manipularea progresiei sau a resurselor economice prin rețele nesigure \textit{Peer-to-Peer} (P2P), proiectul implementează o arhitectură \textit{Server-Authoritative} strictă, validată centralizat.

\section{Obiectivele proiectului}
Pentru a soluționa provocările identificate, lucrarea urmărește trei obiective arhitecturale principale:
\begin{enumerate}
    \item \textbf{Model Server-Authoritative strict:} Izolarea logicii decizionale în instanțe de joc \textit{Godot Headless}, garantând că backend-ul este singurul arbitru al stării jocului și al economiei.
    \item \textbf{Matchmaking rezilient:} Construirea unui algoritm tranzacțional atomic, capabil să proceseze concurent cererile jucătorilor, eliminând coliziunile de memorie (\textit{race conditions}).
    \item \textbf{Sincronizare asincronă non-blocantă:} Proiectarea unui mecanism de \textit{debouncing} pe client care elimină suprasolicitarea rețelei (\textit{spamming} HTTP) și maschează latența.
\end{enumerate}

\textit{Notă privind structura lucrării: Teza este organizată astfel încât să reflecte fluxul decizional tehnic: Capitolul 2 justifică ecosistemul tehnologic ales, Capitolul 3 detaliază arhitectura sistemului distribuit, Capitolul 4 analizează algoritmii concurenți de matchmaking, iar Capitolul 5 validează empiric eficiența și siguranța soluției propuse.}